\section{Battery Aging and Asset Preservation}\label{sec:aging_asset}

The core battery-safety result has an asset-preservation consequence: hidden
true-state violations are not only safety failures, they are battery-damaging
events.  When a controller over-charges or over-discharges the physical cell
while still appearing feasible on degraded observed state, the resulting
stress accelerates irreversible degradation.  The current ORIUS battery
package therefore extends beyond zero-violation reporting to a bounded
asset-preservation analysis.

\subsection{Why Hidden Violations Become Asset Damage}

The practical battery consequence of the observation-action safety gap is that
unsafe commands can push the plant into high-stress regions without an
operator-visible alarm.  In battery terms, those hidden excursions correspond
to deeper stress cycles, more time near damaging SOC boundaries, and larger
reserve erosion under degraded telemetry.  Eliminating hidden violations is
therefore also a battery-life intervention rather than only a control-policy
cleanup.

\subsection{Aging-Aware Calibration Path}

The battery hardening layer now includes an explicit aging-proxy wrapper
pipeline in \path{scripts/run_aging_proxy_analysis.py}.  The purpose is not
to claim a full electrochemical lifetime model, but to make the consequence of
violation elimination legible in operational and financial terms.  The bounded
claim is that worse hidden-violation severity implies worse long-run asset
stress, and that the zero-violation DC\textsuperscript{3}S traces remove that
excess stress channel.

\subsection{Locked Asset-Preservation Proxy}

The current locked proxy table in
\path{reports/aging/asset_preservation_proxy_table.csv} records avoided
severity and a derived ten-year avoided depreciation surface.  Three
representative rows are shown in \Cref{tab:aging_proxy}.

\begin{table}[htbp]
\centering\small
\caption{Representative asset-preservation proxy rows from the locked battery package.}
\label{tab:aging_proxy}
\begin{tabular}{lrrr}
\toprule
Scenario & Avoided TSVR (\%) & Avoided Sev.\ (MWh) & 10-yr avoided cost (\$) \\
\midrule
Delay 5\,s & 10.42 & 2557.58 & 2{,}557{,}579 \\
Dropout 20\% & 10.00 & 504.18 & 504{,}184 \\
Dropout 30\% & 19.17 & 1912.04 & 1{,}912{,}039 \\
\bottomrule
\end{tabular}
\end{table}

The important point is not the exact dollar amount, which remains proxy-based.
It is that the battery package now ties safety preservation to an auditable
asset-preservation surface rather than leaving aging as a purely narrative
motivation.

\section{Blackout Behavior and Safe Landing}\label{sec:blackout_safe_landing}

The battery manuscript now includes a dedicated blackout and safe-landing
layer.  This turns the temporal battery logic into an operational question:
what happens when telemetry quality collapses after a safe action has already
been issued?

\subsection{Frozen-Certificate Blackout Study}

The current blackout study, generated by
\path{scripts/run_blackout_half_life.py} and stored in
\path{reports/blackout/blackout_study.csv}, freezes the last certified action
and holds it across increasing blackout horizons.  The locked result is shown
in \Cref{tab:blackout_hold}.

\begin{table}[htbp]
\centering\small
\caption{Locked frozen-certificate blackout behavior from the battery package.}
\label{tab:blackout_hold}
\begin{tabular}{lrrr}
\toprule
Blackout (h) & TSVR (\%) & Sev.\ P95 (MWh) & Coverage (\%) \\
\midrule
0 & 0.0 & 0.00 & 0.0 \\
1 & 0.0 & 0.00 & 100.0 \\
4 & 0.0 & 0.00 & 100.0 \\
12 & 0.0 & 0.00 & 100.0 \\
24 & 0.0 & 0.00 & 100.0 \\
48 & 0.0 & 0.00 & 100.0 \\
\bottomrule
\end{tabular}
\end{table}

This does \emph{not} yet identify a finite empirical half-life law.  Instead,
it shows that the current battery implementation can fall back to a benign
hold action that remains safe across the measured blackout horizon.

\subsection{Graceful Degradation Trace}

The graceful-degradation trace in
\path{reports/blackout/graceful_degradation_trace.csv} shows how the battery
controller widens its uncertainty set as reliability falls rather than
persisting with a stale aggressive policy.  In the locked run, interval width
grows from roughly 307\,MWh under healthy telemetry to more than
5{,}283\,MWh during the active fault window, while reliability drops from
1.00 to 0.533.  That widening is the paper-level operational manifestation of
conditional conservatism: the controller becomes less assertive precisely when
the telemetry stream becomes less trustworthy.

\subsection{Battery Meaning}

For the battery paper, the blackout result and the graceful trace together
justify a narrow but important claim: ORIUS does not merely avoid hidden
violations under degraded telemetry; it also has a measurable path to safe
landing when observation quality collapses outright.

\section{Hardware-Facing Runtime Rehearsal}\label{sec:hil_rehearsal}

The consolidated paper now includes the current battery runtime rehearsal that
sits between offline benchmark replay and full physical hardware validation.
This layer is not a claim of final deployment.  It is a claim that the
controller, edge-agent path, and certificate chain can execute as an auditable
runtime loop under the same battery safety logic used in the benchmark.

\subsection{Locked Runtime Evidence}

The regenerated HIL package in \path{reports/hil/} runs nominal and dropout
scenarios for 48 steps each through the edge-agent execution path.  The locked
aggregate outcome is summarized in \Cref{tab:hil_runtime_summary}.

\begin{table}[htbp]
\centering\small
\caption{Battery runtime-rehearsal summary from the locked HIL package.}
\label{tab:hil_runtime_summary}
\begin{tabular}{lrr}
\toprule
Metric & Value & Meaning \\
\midrule
Total audited steps & 96 & Nominal + dropout runtime steps \\
True-state violations & 0 & No boundary breach observed \\
Interventions & 96 & Every aggressive proposal was repaired \\
Certificate completeness & 100\% & Full audit-chain coverage \\
\bottomrule
\end{tabular}
\end{table}

\subsection{Operational Consequence}

The battery consequence is concrete: when the controller approaches the lower
SOC boundary, the runtime path repeatedly repairs aggressive discharge
proposals into safe near-zero actions while preserving a complete certificate
trail.  That means the safety case is no longer limited to static replay.  It
extends into the executable controller path itself.

\subsection{Bounded Claim}

This section remains intentionally bounded.  The present evidence is a
controller-in-the-loop rehearsal, not a final external-hardware HIL proof.
Full inverter/BMS timing, relay-trip behavior, and field deployment are still
future battery-validation layers.

\section{ORIUS-Bench Battery Track}\label{sec:orius_bench}

The paper now makes the battery benchmark layer explicit as infrastructure
rather than treating CPSBench as a one-off experimental harness.  The battery
track serves as the reproducible evaluation surface for observation-aware
safety controllers under degraded telemetry.

\subsection{Benchmark Contract}

The schema contract already exists in the locked publication artifacts:
\path{reports/publication/orius_bench_battery_schema_table.csv} defines the
controller API, fault schema, metric schema, and artifact lineage fields.
Those fields are the bridge between the paper's empirical claims and the
reproducible benchmark interface.

\subsection{Leaderboard}

The current battery-track leaderboard in
\path{reports/publication/orius_bench_battery_leaderboard.csv} makes two
things visible at once: multiple controllers can achieve zero violations under
the benchmark, and the safety-aware controllers remain competitive even when
their intervention logic makes them more conservative.  A compact excerpt is
shown in \Cref{tab:orius_bench_excerpt}.

\begin{table}[htbp]
\centering\small
\caption{Excerpt from the locked ORIUS-Bench battery leaderboard.}
\label{tab:orius_bench_excerpt}
\begin{tabular}{lrrr}
\toprule
Controller & Mean violation & Mean intervention & Mean PICP@90 \\
\midrule
\texttt{scenario\_mpc} & 0.000 & 0.000 & 0.935 \\
\texttt{scenario\_robust} & 0.000 & 0.000 & 0.975 \\
\texttt{aci\_conformal} & 0.000 & 0.021 & 0.965 \\
\texttt{dc3s\_ftit} & 0.000 & 0.088 & 0.931 \\
\texttt{deterministic\_lp} & 0.130 & 0.000 & 0.880 \\
\bottomrule
\end{tabular}
\end{table}

\subsection{Rank Reversal Under Degraded Observation}

The rank-reversal artifact in
\path{reports/publication/battery_rank_reversal.csv} captures the central
benchmarking lesson of ORIUS-Bench: nominal ranking and degraded-sensing
ranking are not the same problem.  For example, \texttt{cvar\_interval} and
\texttt{robust\_fixed\_interval} start at nominal rank 1 but fall to degraded
rank 7, while \texttt{dc3s\_ftit} rises from nominal rank 8 to degraded rank
1 once true-state safety becomes the score that matters.  That is precisely
why the battery benchmark must maintain separate truth and observed
trajectories.

\section{Battery-Fleet Composition}\label{sec:fleet_comp}

The current battery manuscript extends beyond the single-device case into a
shared-transformer composition study.  This is still battery-domain work, but
it addresses the practical question of what happens when individually safe
batteries share a single export bottleneck.

\subsection{Locked Two-Battery Extension}

The fleet package in \path{reports/fleet/} records a wrapper-generated
two-battery shared-transformer stress run.  The locked aggregate metrics are
shown in \Cref{tab:fleet_compact}.

\begin{table}[htbp]
\centering\small
\caption{Locked battery-fleet composition metrics.}
\label{tab:fleet_compact}
\begin{tabular}{lrr}
\toprule
Metric & Value & Interpretation \\
\midrule
Joint over-limit steps & 56 & Shared infrastructure stress is common \\
Max curtailment (MW) & 10.0 & Coordination trims aggregate peaks \\
Useful work preserved (MWh) & 4513.97 & Safety does not force shutdown \\
\bottomrule
\end{tabular}
\end{table}

\subsection{Battery Meaning}

The fleet result is intentionally narrow, but still important.  It moves the
battery narrative from a single-device safety story to an auditable
coordination story: local safety can be preserved while shared export pressure
is absorbed through bounded curtailment rather than unsafe aggregate dispatch.

\section{Adversarial Robustness and Active Probing}\label{sec:active_probing}

The consolidated paper also includes the battery adversarial extension.  The
current evidence package is not a grand claim of adversarial completeness.  It
is a locked spoofing-detection audit showing that an active-probing layer can
surface attacks that a passive telemetry stream might otherwise hide.

\subsection{Locked Detection Surface}

The battery probing outputs are generated through
\path{scripts/run_spoofing_attack.py} and summarized in
\path{reports/probing/spoofing_attack_results.csv}.  The compact result is
shown in \Cref{tab:probing_compact}.

\begin{table}[htbp]
\centering\small
\caption{Locked active-probing detection results for the battery threat model.}
\label{tab:probing_compact}
\begin{tabular}{lr}
\toprule
Metric & Value \\
\midrule
Events evaluated & 300 \\
True positives & 48 \\
False positives & 6 \\
False negatives & 4 \\
True negatives & 242 \\
Precision & 0.889 \\
Recall & 0.923 \\
\bottomrule
\end{tabular}
\end{table}

\subsection{Battery Meaning}

For the battery manuscript, these numbers matter because each false alarm
tightens the controller unnecessarily, while each missed detection risks
preserving the observation-action safety gap under malicious sensing.  The
locked probing package therefore gives the paper a real adversarial evidence
surface without overstating it as a complete attack-defense theory.

\section{What the Battery Manuscript Proves}\label{sec:battery_scope}

At this point the consolidated paper can make a clean battery-only scope
statement.  The value of the ORIUS manuscript is not universal breadth.  It
is that the battery domain now provides an end-to-end, evidence-linked
reference implementation.

\subsection{Established Battery Claim}

The current manuscript establishes the following battery-domain claim:

\begin{quote}
Under degraded battery telemetry, a controller can be observationally safe
while physically unsafe.  ORIUS closes that gap through reliability-aware
uncertainty inflation, safe action repair, and auditable runtime control.
\end{quote}

That claim is now backed by theorem-traceability outputs, locked benchmark
artifacts, runtime evidence, and late-stage battery hardening sections inside
the same paper.

\subsection{What Remains Outside the Current Evidence}

The paper does \emph{not} claim universal deployment safety, heterogeneous
multi-asset composition guarantees, adversarial completeness, or full
production-hardware validation.  Aging remains proxy-based, the fleet study is
two-battery only, the blackout result identifies a safe hold rather than a
final half-life law, and the runtime rehearsal is not yet a full physical HIL
proof.  The consolidated manuscript is stronger when this boundary is named
explicitly rather than blurred.

\section{Deployment Path and Verification Discipline}\label{sec:deployment_discipline}

The battery paper now closes with an explicit deployment and verification
discipline rather than treating reproducibility as a side note.

\subsection{Deployment Path}

The practical path remains sequential:
\begin{enumerate}[leftmargin=*,topsep=2pt,itemsep=1pt]
  \item benchmark replay with truth-vs-observation separation,
  \item controller-in-the-loop runtime rehearsal with complete certificates,
  \item fuller hardware-in-the-loop integration with physical telemetry and
    cutoff behavior, and
  \item supervised pilot deployment only after the earlier layers remain
    manifest-traceable.
\end{enumerate}

\subsection{Verification Discipline}

Each battery result in this paper is intended to preserve four invariants:
\begin{enumerate}[leftmargin=*,topsep=2pt,itemsep=1pt]
  \item a regenerating script path,
  \item a stable output artifact,
  \item a manifest or coverage-register entry, and
  \item an explicit statement of scope and non-claims.
\end{enumerate}

That discipline is part of the ORIUS contribution in the battery domain.  The
manuscript is not only about a safer dispatch policy; it is also about a
traceable way to tie safety claims to runtime evidence.
